\documentclass[11pt]{article}
\usepackage[margin=1in]{geometry}
\usepackage{amsmath,amssymb,amsthm,mathtools}
\usepackage{enumitem}
\usepackage{microtype}
\usepackage{hyperref}
\usepackage[T1]{fontenc}
\usepackage{lmodern}

\newtheorem{conjecture}{Conjecture}
\newtheorem{lemma}{Lemma}
\newtheorem{question}{Question}
\newcommand{\OP}{\mathcal O}
\newcommand{\Lin}{\operatorname{LinExt}}
\newcommand{\Aut}{\operatorname{Aut}}
\newcommand{\Cyc}{\operatorname{CycExt}}
\newcommand{\Lie}{\operatorname{Lie}}
\newcommand{\sgn}{\operatorname{sgn}}
\newcommand{\Ind}{\operatorname{Ind}}
\newcommand{\ch}{\operatorname{ch}}
\newcommand{\HH}{\mathbb H}

\title{A Poset--Joyal--Klyachko--Stanley Conjecture\\and a Possible Poset Lie Theory}
\author{}
\date{Research note --- October 2026}

\begin{document}
\maketitle

\begin{abstract}
This note records a conjectural equivariant refinement of the theorem of
Jen\v{c}a--Sarkoci on the topology of the lattice of order-preserving
partitions of a finite poset. The conjecture identifies the top homology
representation with a Klyachko/free-Lie projection of the linear-extension
set. It is true for the antichain by the classical
Joyal--Klyachko--Stanley correspondence, has the correct dimension for every
poset by Jen\v{c}a--Sarkoci, and has been checked computationally for every
unlabeled poset through $n=5$, together with several symmetric examples at
$n=6$. The note isolates an equivalent character formula and a concrete
restriction lemma whose proof would settle the conjecture. It also records
a global $S_n$-module obtained by summing over all labeled posets and the
resulting Witt/plethystic-logarithm pattern.
\end{abstract}

\section{The basic objects}

Let $P$ be a finite poset on the labeled ground set $[n]$.

\subsection{Order-preserving partitions}
An \emph{order-preserving partition} (order congruence) of $P$ is a set
partition $\pi$ of $[n]$ such that the quotient relation $P/\pi$ is again
acyclic. Let $\OP(P)$ be the lattice of order-preserving partitions and
write
\[
\widehat{\OP(P)}=\OP(P)\setminus\{\hat0,\hat1\}.
\]

Equivalently, begin with the compatible ordered-partition complex
\[
\mathcal F(P)=\coprod_k \operatorname{Sur}_{\rm iso}(P,[k]),
\]
whose top-dimensional chambers are the linear extensions of $P$.
Forgetting block order gives the support map
\[
\mathcal F(P)\longrightarrow \OP(P).
\]
Geometrically, $P$ cuts out the order cone
\[
C_P=\{x\in\mathbb R^n:x_i\le x_j \text{ whenever }i<_Pj\},
\]
and the braid arrangement subdivides $C_P$ into chambers indexed by
$\Lin(P)$.

\subsection{Cyclic classes of linear extensions}
Let $\Lin(P)\subseteq S_n$ be the set of linear extensions, written as words
$w=w_1\cdots w_n$. Let $C_n=\langle\gamma\rangle$ with
$\gamma=(1\,2\,\cdots\,n)$ act on positions by cyclic rotation. A cyclic
class of linear extensions is a $C_n$-orbit in $S_n$ meeting $\Lin(P)$.
Let
\[
e_C(P)=|\Cyc(P)|.
\]

Jen\v{c}a--Sarkoci prove
\[
\Delta(\widehat{\OP(P)})\simeq\bigvee^{e_C(P)}S^{n-3}.
\]
Hence
\begin{equation}
\dim \widetilde H_{n-3}(\widehat{\OP(P)};\mathbb C)=e_C(P).
\label{eq:JSdim}
\end{equation}
If $P$ is connected, then $e_C(P)=|\Lin(P)|$.

\section{The candidate poset Lie module}
Fix a primitive $n$th root of unity $\zeta$ and set
\[
e_\zeta=\frac1n\sum_{j=0}^{n-1}\zeta^{-j}\gamma^j\in\mathbb C[S_n].
\]
Define
\begin{equation}
\mathsf{Lie}(P):=\operatorname{span}_{\mathbb C}\{we_\zeta:w\in\Lin(P)\}
\subseteq\mathbb C[S_n]e_\zeta.
\label{eq:posetLie}
\end{equation}

Two linear extensions in the same cyclic class determine proportional
vectors after multiplication by $e_\zeta$, while vectors supported on
different right $C_n$-cosets are independent. Therefore
\begin{equation}
\dim \mathsf{Lie}(P)=e_C(P).
\label{eq:posetLiedim}
\end{equation}
The automorphism group $G=\Aut(P)$ acts on $\mathsf{Lie}(P)$ by left
multiplication.

\begin{conjecture}[Poset JKS conjecture]
For every finite $n$-element poset $P$, there is a natural
$\Aut(P)$-equivariant isomorphism
\begin{equation}
\boxed{
\widetilde H_{n-3}(\widehat{\OP(P)};\mathbb C)
\cong \sgn\otimes\mathsf{Lie}(P).}
\label{eq:mainconj}
\end{equation}
Here $\sgn$ is the restriction to $\Aut(P)$ of the sign character of $S_n$.
\end{conjecture}

\section{Why this is the correct generalization}

\subsection{The antichain recovers classical free Lie theory}
If $P$ is the antichain on $[n]$, then
\[
\OP(P)=\Pi_n,\qquad \Lin(P)=S_n.
\]
Thus
\[
\mathsf{Lie}(P)=\mathbb C[S_n]e_\zeta\cong\Ind_{C_n}^{S_n}\zeta.
\]
By Klyachko this is the multilinear degree-$n$ component $\Lie_n$ of the
free Lie algebra. The conjecture becomes
\[
\widetilde H_{n-3}(\overline{\Pi_n})\cong\sgn\otimes\Lie_n,
\]
the classical Joyal--Klyachko--Stanley theorem.

Thus the conjecture is not merely analogous to JKS: it is a proposed
restriction of the same phenomenon from the full braid arrangement to the
order cone of a poset.

\subsection{Compatible faces and the braid arrangement}
The braid arrangement has
\[
\text{chambers}=\text{permutations},\qquad
\text{faces}=\text{ordered set partitions},
\]
and support lattice $\Pi_n$. Restriction to $C_P$ retains exactly the
chambers $\Lin(P)$; compatible faces are surjective isotone maps $P\to[k]$,
and their supports are the elements of $\OP(P)$.

\section{Equivalent character formulation}
Let
\[
H_P=\widetilde H_{n-3}(\widehat{\OP(P)};\mathbb C).
\]
For $g\in G=\Aut(P)$, homology is concentrated in top degree, so the
Lefschetz trace formula gives
\begin{equation}
\chi_{H_P}(g)=(-1)^{n-3}\mu(\OP(P)^g),
\label{eq:homchar}
\end{equation}
where $\OP(P)^g$ is the fixed-point subposet and $\mu$ its M\"obius number.

A cyclic class $C\in\Cyc(P)$ contributes to the trace on $\mathsf{Lie}(P)$
only if it is setwise fixed. For a fixed class, write
\[
r_g(C)\in\mathbb Z/n
\]
for the rotation amount induced by $g$. Then
\begin{equation}
\chi_{\mathsf{Lie}(P)}(g)=
\sum_{\substack{C\in\Cyc(P)\\gC=C}}\zeta^{r_g(C)}.
\label{eq:liechar}
\end{equation}
Hence the conjecture is equivalent to
\begin{equation}
\boxed{
(-1)^{n-3}\mu(\OP(P)^g)
=\sgn(g)\sum_{\substack{C\in\Cyc(P)\\gC=C}}\zeta^{r_g(C)}}
\qquad(g\in\Aut(P)).
\label{eq:characterconj}
\end{equation}
At $g=1$ this reduces exactly to the Jen\v{c}a--Sarkoci counting theorem.
Thus the conjecture is an equivariant root-of-unity refinement of their
result.

\section{Embedding into the classical JKS module}
There is an inclusion
\[
\OP(P)\hookrightarrow\Pi_n.
\]
Consequently
\[
\Delta(\widehat{\OP(P)})\subseteq\Delta(\overline{\Pi_n}).
\]
Both complexes have top dimension $n-3$. There are no ambient
$(n-2)$-chains, so a nonzero top cycle in $\widehat{\OP(P)}$ cannot become a
boundary after inclusion. Thus
\begin{equation}
H_P\hookrightarrow\widetilde H_{n-3}(\overline{\Pi_n})
\label{eq:homembed}
\end{equation}
is injective and $\Aut(P)$-equivariant.

Let
\[
J:\widetilde H_{n-3}(\overline{\Pi_n})\otimes\sgn
\stackrel{\sim}{\longrightarrow}\Lie_n
\]
be a classical JKS isomorphism, and let
\[
W_P=\mathbb C\{\Lin(P)\}\subseteq\mathbb C[S_n].
\]
The conjecture follows from the single restricted-image statement
\begin{equation}
\boxed{
J(H_P\otimes\sgn)
=\operatorname{span}\{we_\zeta:w\in\Lin(P)\}.}
\label{eq:restrictionlemma}
\end{equation}
This is the main concrete lemma still to be proved.

\section{A second formulation using Lie idempotents}
One may replace $e_\zeta$ by a standard Lie idempotent, for example a
Klyachko idempotent $\kappa_n$, since both generate isomorphic copies of the
Lie representation in $\mathbb C[S_n]$. The expected statement is
\[
J(H_P\otimes\sgn)=W_P\kappa_n\cong W_Pe_\zeta.
\]
This suggests a group-algebra proof:
\begin{enumerate}[label=\arabic*.]
\item choose an explicit combinatorial JKS map $J$;
\item evaluate it on generators of the top cycles of $\OP(P)$;
\item prove the image lies in $W_P\kappa_n$;
\item use $\dim H_P=e_C(P)=\dim W_Pe_\zeta$ to conclude equality.
\end{enumerate}
Only containment is genuinely new.

\section{Potential proof via splitting spheres}
For a linear extension $w=w_1\cdots w_n$, allow cuts only between successive
letters of $w$. These interval partitions form a Boolean sublattice whose
proper part is an $(n-3)$-sphere. Classical proofs of JKS, especially those
of Barcelo and Wachs, use closely related splitting cycles/bracketings to
compare partition-lattice homology with free Lie algebra bases.

This is promising for \eqref{eq:restrictionlemma}. A caution is essential:
individual cyclically rotated splitting cycles need not be scalar multiples
in homology, whereas
\[
w\gamma^je_\zeta=\zeta^jwe_\zeta
\]
(up to Fourier-sign convention). Thus the cyclic Fourier projection must be
introduced after assembling the relevant top cycles; one should not expect a
sphere-by-sphere identification.

\begin{lemma}[Desired restricted JKS lemma]
Under a suitable explicit JKS isomorphism, the span of top cycles supported
on order-preserving partitions of $P$ maps onto $W_Pe_\zeta$.
\end{lemma}

\section{Potential proof via hyperplane-arrangement Lie theory}
Aguiar--Mahajan develop a Lie theory relative to finite real hyperplane
arrangements, extending the classical JKS correspondence beyond the braid
arrangement. Their framework uses the Tits algebra, faces, flats, lunes and
arrangement-theoretic Lie elements.

The poset setting is a cone restriction of the braid arrangement rather
than a new central arrangement. The natural program is:
\begin{enumerate}[label=\arabic*.]
\item express $\mathcal F(P)$ as the face/left-regular-band object attached
to the cone $C_P$;
\item identify its support semilattice with $\OP(P)$;
\item identify the maximal-support Lie/Peirce component;
\item show that inclusion into the full braid-arrangement Tits algebra sends
that component to $W_Pe_\zeta\subseteq\Lie_n$.
\end{enumerate}
If the existing theory applies directly to cones/lunes, this may prove the
conjecture conceptually; otherwise it isolates the needed cone-restriction
version of JKS.

\section{Computational evidence}
The full character identity \eqref{eq:characterconj} has been checked
computationally for every unlabeled poset on $n\le5$. In particular, for
$n=3,4,5$ this covers $5+16+63=84$ nontrivial isomorphism types.

For each $P$ and each $g\in\Aut(P)$ the computation compared
\[
(-1)^{n-3}\mu(\OP(P)^g)
\]
with
\[
\sgn(g)\sum_{C:gC=C}\zeta^{r_g(C)}.
\]
No discrepancy was found.

Additional symmetric $n=6$ examples tested include a five-armed star, the
complete bipartite $3+3$ poset, and three disjoint $2$-chains. In the last
example an automorphism cyclically permuting the three components gives
character value $-1$, so the test detects a genuine nontrivial root-of-unity
phase rather than merely fixed-point counting.

These computations should be independently reproduced and archived before
being cited as formal results.

\section{A global $S_n$-module}
Restore full symmetric-group symmetry by summing over every labeled poset on
$[n]$:
\begin{equation}
\HH_n:=\bigoplus_{P\text{ on }[n]}
\widetilde H_{n-3}(\widehat{\OP(P)};\mathbb C).
\label{eq:globalH}
\end{equation}
Relabeling permutes the summands, so $\HH_n$ is an $S_n$-module.
Conjecturally,
\[
\HH_n\otimes\sgn\cong\bigoplus_{P\text{ on }[n]}\mathsf{Lie}(P).
\]

Fix one cyclic order and let $B_n$ be the set of labeled posets compatible
with it. The expected induced representation is
\begin{equation}
\HH_n\cong\sgn\otimes
\Ind_{C_n}^{S_n}(\zeta\otimes\mathbb C[B_n]).
\label{eq:globalind}
\end{equation}
This places the global object in the theory of cyclic/Foulkes characters.

\section{Witt and plethystic-logarithm pattern}
Let $a_n$ be the number of naturally labeled posets on $[n]$ and set
\[
A(x)=1+\sum_{n\ge1}a_nx^n.
\]
The initial values are
\[
a_n=1,2,7,40,357,4824,96428,\ldots.
\]
Empirically, for $b_n=|B_n|$,
\begin{equation}
b_n=n[x^n]\log A(x),
\label{eq:bnlog}
\end{equation}
which gives
\[
b_n=1,3,16,131,1561,26532,637890,21580043,\ldots.
\]

Assuming the global induced-module description, the Frobenius characteristic
has the divisor-sum form
\begin{equation}
\ch\HH_n=
\frac1n\sum_{d\mid n}(-1)^{n-n/d}\mu(d)b_{n/d}p_d^{n/d},
\label{eq:frob}
\end{equation}
with the displayed sign depending on the convention for the outer sign
twist. Ignoring that twist,
\begin{equation}
\sum_{n\ge1}\ch(\HH_n^{\rm untwisted})
=\sum_{d\ge1}\frac{\mu(d)}{d}\log A(p_d),
\label{eq:plethlog}
\end{equation}
which has exactly the formal shape of a plethystic logarithm.

For comparison,
\[
\ch(\Lie_n)=\frac1n\sum_{d\mid n}\mu(d)p_d^{n/d}.
\]
Thus the global poset construction appears to replace the scalar one-generator
input in Witt's formula by the naturally labeled-poset series $A$.

\section{Prime-degree sanity check}
For prime $p$, the $C_p$-action on $B_p$ has the antichain as a fixed point
and every other orbit has size $p$. The induced-module picture predicts
\[
\HH_p\cong\frac{b_p-1}{p}\,\mathbb C[S_p]\oplus
(\sgn\otimes\Lie_p).
\]
For example,
\[
\HH_3\cong5\,\operatorname{Reg}_{S_3}\oplus
(\sgn\otimes\Lie_3),
\]
\[
\HH_5\cong312\,\operatorname{Reg}_{S_5}\oplus
(\sgn\otimes\Lie_5).
\]
The nonregular summand is exactly the classical Lie representation.

\section{Status: proved, supported, and open}
\subsection*{Known or elementary}
\begin{itemize}
\item Jen\v{c}a--Sarkoci: $\Delta(\widehat{\OP(P)})$ is a wedge of
$e_C(P)$ spheres of dimension $n-3$.
\item Hence $\dim H_P=e_C(P)$.
\item By construction $\dim\mathsf{Lie}(P)=e_C(P)$.
\item The antichain case is classical JKS/Klyachko.
\item The inclusion $H_P\hookrightarrow\widetilde
H_{n-3}(\overline{\Pi_n})$ is equivariant and injective.
\item The main conjecture is equivalent to the character identity
\eqref{eq:characterconj}.
\end{itemize}

\subsection*{Computationally supported}
\begin{itemize}
\item The full character identity holds for every unlabeled poset through
$n=5$ in the current computations.
\item It also holds for several symmetric $n=6$ examples.
\item The global dimensions and cyclic orbit structure are consistent with
\eqref{eq:globalind}--\eqref{eq:plethlog}.
\end{itemize}

\subsection*{Still to prove}
\begin{itemize}
\item The restricted JKS image identity \eqref{eq:restrictionlemma}.
\item Equivalently, the fixed-point M\"obius/root-of-unity identity
\eqref{eq:characterconj}.
\item Naturality with respect to useful operations on posets.
\item The global induced-module formula \eqref{eq:globalind}.
\item A direct proof of \eqref{eq:bnlog} in the precise convention used for
$B_n$.
\item Whether the modules $\mathsf{Lie}(P)$ admit an intrinsic bracket,
operadic composition, or Hopf-monoid structure deserving the name ``poset
Lie algebra.''
\end{itemize}

\section{Concrete next research steps}
\subsection{1. Prove the character identity directly}
For fixed $g\in\Aut(P)$, classify $\OP(P)^g$ and seek a $g$-equivariant
version of the Jen\v{c}a--Sarkoci recursion. The target is
\[
(-1)^{n-3}\mu(\OP(P)^g)=
\sgn(g)\sum_{C:gC=C}\zeta^{r_g(C)}.
\]
The new issue is that fixed cyclic classes must be counted with phase, not
merely multiplicity.

\subsection{2. Prove restricted JKS containment}
Choose an explicit JKS map from the literature, preferably one using
splitting cycles or bracketings. For every top cycle supported in $\OP(P)$,
show
\[
J(z)\in W_Pe_\zeta.
\]
Dimension equality then gives the theorem.

\subsection{3. Compare with arrangement Lie idempotents}
Translate $C_P$ and $\mathcal F(P)$ into the Aguiar--Mahajan framework and
check whether the maximal-support Lie/Peirce component is already known to
be the restriction of the braid-arrangement Lie component.

\subsection{4. Extend computation strategically}
Rather than enumerate all posets at $n=6,7$, prioritize highly symmetric
families: unions of equal chains, complete layered posets, crowns, products,
and symmetric width-$2$/$3$ examples. Record for every conjugacy class
\[
\mu(\OP(P)^g),\quad \#\{C:gC=C\},\quad\{r_g(C)\},\quad\chi_{H_P}(g).
\]

\subsection{5. Investigate functorial composition}
Test $\mathsf{Lie}(P)$ under disjoint sum, ordinal sum, products, and poset
substitution. A genuine species/operad structure should explain the
plethystic logarithm rather than merely reproduce it numerically.

\section{Relation to the moonshine motivation}
The construction does not presently produce modular forms. What it does
produce is a structurally analogous intermediate phenomenon:
\[
\text{large combinatorial family}
\longrightarrow\text{equivariant homology}
\longrightarrow\text{Lie/Foulkes characters}
\longrightarrow\text{Witt/plethystic-logarithm identities}.
\]
This is notable because denominator/Witt identities are also a central
algebraic layer in Borcherds-style moonshine constructions. The immediate
question is therefore not whether this is already moonshine, but whether the
compatible-face/order-congruence construction defines a natural
Lie-theoretic primitive functor on posets.

\begin{question}
Does the compatible-face/order-congruence construction define a natural
primitive or Lie functor on posets whose local character theory is governed
by \eqref{eq:mainconj} and whose global cycle index is the plethystic
logarithm of the poset species?
\end{question}

A positive answer would provide a coherent poset Lie theory. Only then does
it become natural to search for an additional grading or operator whose
traces might have modular or automorphic properties.

\section{References}
\begin{thebibliography}{99}

\bibitem{JencaSarkoci}
G.~Jen\v{c}a and P.~Sarkoci,
\emph{Linear extensions and order-preserving poset partitions},
Journal of Combinatorial Theory, Series A \textbf{122} (2014), 28--38.
DOI: \href{https://doi.org/10.1016/j.jcta.2013.09.010}
{10.1016/j.jcta.2013.09.010}.

\bibitem{Barcelo}
H.~Barcelo,
\emph{On the action of the symmetric group on the free Lie algebra and the
partition lattice},
Journal of Combinatorial Theory, Series A \textbf{55} (1990), 93--129.
DOI: \href{https://doi.org/10.1016/0097-3165(90)90050-7}
{10.1016/0097-3165(90)90050-7}.

\bibitem{Wachs}
M.~L.~Wachs,
\emph{On the (co)homology of the partition lattice and the free Lie algebra},
Discrete Mathematics \textbf{193} (1998), 287--319.
DOI: \href{https://doi.org/10.1016/S0012-365X(98)00147-2}
{10.1016/S0012-365X(98)00147-2}.

\bibitem{AguiarMahajan}
M.~Aguiar and S.~Mahajan,
\emph{Bimonoids for Hyperplane Arrangements},
Encyclopedia of Mathematics and its Applications, Cambridge University
Press, 2020.

\bibitem{Klyachko}
A.~A.~Klyachko,
\emph{Lie elements in the tensor algebra},
Siberian Mathematical Journal \textbf{15} (1974).
The relevant theorem identifies the multilinear Lie representation with the
representation induced from a faithful character of a long cyclic subgroup;
exact pagination should be checked before publication.

\end{thebibliography}

\section*{Research status summary}
The compact statement to carry forward is
\[
\boxed{
\widetilde H_{n-3}(\widehat{\OP(P)})
\stackrel{?}{\cong}
\sgn\otimes\operatorname{span}\{we_\zeta:w\in\Lin(P)\}.}
\]
The strongest next move is to prove either the restricted JKS image identity
\eqref{eq:restrictionlemma} or the equivalent fixed-point M\"obius identity
\eqref{eq:characterconj}.

\end{document}
